% =====================================================================
%  CONCLUSIONES (versión por objetivos)
%
%  Parte de las conclusiones que aprobaron los autores (commit 8fa9f2a). El 11 de
%  noviembre de 2026, por corrección del director, cada párrafo empieza por lo que
%  se hizo para su objetivo y se quitaron las causas, que pasaron a la discusión:
%  la planificación sin ROS 2, la falta de codificadores, la banda muerta como
%  causa de la tolerancia, la tarjeta de dos núcleos y lo que la instrumentación
%  habría registrado sin la confirmación. Lo demás conserva su redacción.
% =====================================================================
\chapter*{Conclusiones}
\addcontentsline{toc}{chapter}{CONCLUSIONES}

Al modelar la arquitectura funcional, con sus requerimientos, sus escenarios de
operación en los edificios Gregorio y Catalayud, una asignación de tareas por piso y
un esquema de comunicación centrado en un coordinador, se comprobó que el cambio de
piso puede resolverse como un problema de coordinación, sin que ningún agente suba
escaleras ni use el ascensor. En ninguna misión de las dos campañas falló la
asignación ni el relevo. Con un solo agente admisible por piso, la asignación resulta
trivial; la dificultad y el aporte del trabajo están en coordinar dos agendas que
dependen una de otra durante el relevo.

Se desarrolló la plataforma robótica cargando el sistema en dos agentes, que integran
locomoción, sensado, procesamiento y comunicación. La plataforma AWS
DeepRacer sirve como agente de guiado en interiores, pero su desempeño depende del
sitio y de la tracción. En el edificio Gregorio la odometría registró como máximo el
73~\% del avance real, mientras que en Catalayud, con el mismo estimador, cumplió el
criterio de la prueba. El sitio de operación debe elegirse como parte del diseño. Por
debajo de 0,40~m/s la tracción no responde, y la tolerancia de llegada se fijó en
0,5~m, el doble de la admitida en simulación. La coordinación se llevó a un computador
portátil, y la duración de las sesiones queda limitada por las baterías, que en las
pruebas se agotaron en los agentes y en ese computador.

Se programó la interfaz móvil como una página web servida en la red local, y para
operar el sistema basta con ella: en las pruebas con los agentes físicos cada misión
se pidió, se siguió y se confirmó desde un teléfono, sin instalar nada y sin salida a
internet. Sin embargo, las funciones que más la acercan a una operación real no
estaban en el anteproyecto y se añadieron al ver el sistema en uso. Con la
confirmación de cambio de piso, una misión no se da por completada con el usuario en
el piso equivocado, y la cancelación devolvió al agente a su escalera en las seis
misiones canceladas. Por lo tanto, los requisitos de una interfaz como esta no pueden
cerrarse de antemano.

Con la evaluación experimental, hecha con un protocolo definido antes de instrumentar
el sistema, se puede afirmar, con 95~\% de confianza, que la tasa de éxito del sistema
supera el 70~\%, tanto en las treinta misiones en simulación como en las nueve con los
agentes físicos, aunque esas muestras no permiten precisar cuánto la supera. Con los
agentes físicos, el tiempo total mediano de una misión fue de 87,6~s. Todos los fallos
fueron de navegación, sobre todo de precisión de llegada, y ninguno de asignación ni
de relevo, de modo que la mejora del sistema pasa primero por la navegación.
